%\documentclass[12pt,oneside]{amsart}
\documentclass{scrartcl}

%\documentclass{article}

\usepackage{amsthm,amsmath,amssymb}
\usepackage[numbers]{natbib}

%\usepackage[colorlinks]{hyperref}
\usepackage{hyperref}
\hypersetup{
    colorlinks,%
    citecolor=black,%
    filecolor=black,%
    linkcolor=black,%
    urlcolor=black
}
\usepackage{hypernat}
\usepackage[top=2.5cm, bottom=2.5cm, left=2.8cm,
right=2.8cm]{geometry}
\usepackage{graphicx}


%\setlength{\parskip}{0pt}
%\setlength{\parsep}{0pt}
%\setlength{\headsep}{0pt}
%\setlength{\topskip}{0pt}
%\setlength{\topmargin}{0pt}
%\setlength{\topsep}{0pt}
%\setlength{\partopsep}{0pt}

\linespread{1}

%\usepackage{marvosym}

%\textwidth = 6.5 in \textheight = 9.2 in \oddsidemargin = 0.0 in
%\evensidemargin = 0.0 in \topmargin = -0.0 in \headheight = 0.0 in
%\headsep = 0.0in \parskip = 0.0in \parindent = 0.0in
%\def\baselinestretch{0.871}


\newtheorem{theorem}{Theorem}[section]
\newtheorem{proposition}[theorem]{Proposition}
\newtheorem{problem}[theorem]{Problem}
\newtheorem{fact}[theorem]{Fact}
\newtheorem{lemma}[theorem]{Lemma}
\newtheorem{defn}[theorem]{Definition}
\newtheorem{corollary}[theorem]{Corollary}
\newtheorem{remark}{Remark}[section]
\newtheorem{example}[theorem]{Example}

\newcommand{\E}{{\mathbb E}}
\newcommand{\se}{{\mathcal{E}}}
\newcommand{\R}{{\mathbb R}}
\renewcommand{\P}{{\mathbb P}}
\newcommand{\ac}{{\mathcal{A}}}
\newcommand{\A}{{\mathcal{A}}}
\newcommand{\B}{{\mathcal{B}}}
\newcommand{\C}{{\mathcal{C}}}
\newcommand{\M}{{\mathcal{M}}}
\newcommand{\Mf}{{\mathfrak{M}}}
\newcommand{\T}{{\mathcal{T}}}
\newcommand{\Z}{{\mathcal{Z}}}
\newcommand{\K}{{\mathcal{K}}}
\newcommand{\lc}{{\mathcal{L}}}
\newcommand{\Ps}{{\mathcal{P}}}
\newcommand{\G}{{\mathcal{G}}}
\newcommand{\pa}{{\mathring{p}}}
\newcommand{\F}{{\cal F}}
\newcommand{\samp}{{\mathcal{X}}}
\newcommand{\X}{{\mathcal{X}}}
\newcommand{\hi}{{\hat{\Phi}}}
\newcommand{\data}{{\text{data}}}
\newcommand{\relu}{{\text{ReLU}}}

\newcommand{\ridge}{{\hat{\beta}_{\text{ridge}}(\lambda)}}
\newcommand{\lasso}{{\hat{\beta}_{\text{lasso}}(\lambda)}}
\newcommand{\ridgemu}{{\hat{\mu}_t^{\text{ridge}}(\lambda)}}
\newcommand{\lassomu}{{\hat{\mu}_t^{\text{lasso}}(\lambda)}}


\newcounter{rcnt}[section]
\renewcommand{\thercnt}{(\roman{rcnt})}

\def\qt#1{\qquad\text{#1}}

\def\argmin{\mathop{\rm argmin}}
\def\argmax{\mathop{\rm argmax}}


\setlength{\parskip}{1.4 \medskipamount}
%\sloppy
%\linespread{1.3}

\begin{document}

\title{STAT 238 - Bayesian Statistics  \\
  Lecture Twenty Four} 
\subtitle{Spring 2026, UC Berkeley}

\author{Aditya Guntuboyina}


\date{30 March 2026}

\maketitle

Our next topic is Bayesian computation. The two main classes of
techniques are Markov Chain Monte Carlo (MCMC) and Variational
Inference (VI). We will start discussing MCMC today. 

\section{Markov Chain Monte Carlo}

The goal is to obtain samples from a target distribution with density 
$\pi$ (in our applications, $\pi$ will be the posterior density). We
assume that $\pi$ is known only upto a normalizing constant. The
normalizing constant can be determined by integrating the unnormalized
$\pi$ (over its domain) but this integration is assumed to be
difficult. 

We
want to generate $\theta^{(0)}, \theta^{(1)}, \dots, \theta^{(N)}$
such that  
\begin{equation}\label{appr}
  \frac{1}{N} \sum_{t=1}^N g(\theta^{(t)}) \approx \int g(\theta)
  \pi(\theta) d\theta. 
\end{equation}
for all functions $g$.

$\theta^{(t)}$ is generated sequentially or iteratively for $t = 0, 1,
2, \dots$ with 
$\theta^{(t)}$ (and some external randomness) used to generate
$\theta^{(t+1)}$.

The simplest MCMC algorithm is the Random Walk Metropolis Algorithm. 

\section{The Random Walk Metropolis Algorithm}

The random-walk Metropolis algorithm works by
generating proposals via a random walk-step from the current state and
then accepting or rejecting them with certain probability. Suppose
that the current state is $\theta^{(t)}$. It then proposes to move to
the state $\theta^{(t)} + Z$ where $Z$ is a random vector having a
symmetric density around 0. Note that the density 
of the proposal $\theta^{(t)} + Z$ is given by $q(\theta^{(t)},
\cdot)$ where 
\begin{equation*}
  q(x, y) = f_Z(y - x). 
\end{equation*}
with $f_Z$ denoting the density of $Z$. Because $f_Z$ is symmetric
about 0, we have  
\begin{equation}\label{symm}
  q(x, y) = q(y, x). 
\end{equation}
The function $q(\cdot, \cdot)$ is commonly referred to as the
candidate-generating density or proposal-generating density. We say
that $q$ is symmetric if \eqref{symm} holds. Clearly random walk
Metropolis with a symmetric $f_{Z}$ uses symmetric proposal-generating
densities. With this notation, the random-walk Metropolis algorithm can be
rephrased as follows. For $x$ and $y$, let  
\begin{equation}\label{acc.symm}
  \alpha(x, y) := \min \left(\frac{\pi(y)}{\pi(x)}, 1\right). 
\end{equation}
$\alpha(x, y)$ is referred to as the \textbf{acceptance probability}
while going from $x$ to $y$. 

\begin{enumerate}
\item Initialize with an arbitrary value $\theta^{(0)}$.
\item Repeat the following for $t = 1, \dots, N$:
  \begin{enumerate}
  \item Let $x = \theta^{(t)}$. 
  \item Generate a candidate or proposal $y$ from the
    \textbf{symmetric} candidate generating density $q(x,
    \cdot)$ and a uniform random number $u \sim \text{Unif}(0,
    1)$.
  \item If $u \geq \alpha(x, y)$, then set
    $\theta^{(t + 1)} = x$. We say in this case that the
    proposed move from $x$ to $y$ is rejected and we stay at $x$. 
  \item If $u < \alpha(x, y)$, then set
    $\theta^{(t+1)} = y$. We say in this case that the
    proposed move from $x$ to $y$ is accepted. 
  \end{enumerate}
\item Return the values $\theta^{(1)}, \dots, \theta^{(N)}$.   
\end{enumerate}
The above algorithm is known as the \textbf{Metropolis Algorithm}. It
works for any symmetric (i.e., satisfying \eqref{symm})
proposal-generating density $q(\cdot, \cdot)$. The Random-Walk
Metropolis algorithm is a special case of this where $q(x,
y)$ is of the form $f_{Z}(x - y)$ for a symmetric (about 0)
density $f_{Z}$.

The Metropolis algorithm can be further generalized to deal with
non-symmetric proposal generating densities. This generalization is
referred to as the Metropolis-Hastings algorithm. We shall look at
this algorithm in the next lecture. We will try to answer the
following questions in the coming lectures: 
\begin{enumerate}
\item What is the intuition behind the form of these algorithms?
\item Why do they actually work and give samples satisfying
  \eqref{appr}? 
\end{enumerate}

To answer these questions, we need to know a little bit about Markov
Chains. The algorithms above output samples $\theta^{(0)}, \dots,
\theta^{(N)}$. Clearly, $\theta^{(t+1)}$ is generated using only the value
$\theta^{(t)}$. In other words, the values $\theta^{(0)}, \dots,
\theta^{(t-1)}$ are not used at all in order to generate
$\theta^{(t+1)}$. Thus $\theta^{(0)}, \dots, \theta^{(N)}$ can be seen
as a realization of a sequence of random vectors $\Theta^{(0)}, \dots,
\Theta^{(N)}$ which form a Markov Chain. Properties of Markov Chains
are used to justify \eqref{appr}, as we shall see in the coming
lectures. 





%\begin{thebibliography}{99}

%\bibitem[Chib and Greenberg(1995)]{chib1995}
%Chib, S. and Greenberg, E. (1995).
%Understanding the Metropolis--Hastings algorithm.
%\textit{The American Statistician}, \textbf{49}, 327--335.


%\bibitem[MacKay and Peto(1995)]{MacKay1995HierarchicalDirichlet}
%D.~J.~C. MacKay and L.~C. Peto,
%\newblock ``A hierarchical Dirichlet language model,''
%\newblock \emph{Natural Language Engineering}, vol.~1,
%no.~3, pp.~289--308, 1995.

%\bibitem[Rubin(1981)]{Rubin1981BayesianBootstrap}
%D.~B. Rubin,
%\newblock ``The Bayesian bootstrap,''
%\newblock \emph{The Annals of Statistics}, vol.~9,
%no.~1, pp.~130--134, 1981.
  
%\bibitem[Fisher(1934)]{Fisher1934}
%R.~A. Fisher,
%\newblock ``Probability, likelihood and quantity of information in the logic of
%uncertain inference,''
%\newblock \emph{Proceedings of the Royal Society of London. Series~A,
%Containing Papers of a Mathematical and Physical Character}, vol.~146,
%no.~856, pp.~1--8, 1934.

%\bibitem[Efron(2010)]{Efron}
%Efron, B. (2010)
%\textit{Large-scale inference: empirical Bayes methods for estimation,
%testing, and prediction}. Cambridge University Press, 2010.

%\bibitem[Shumway and Stoffer(2017)]{ShumwayStoffer}
%Shumway, R. H. and Stoffer, D. S. (2017)
%\textit{Time Series Analysis and Its Applications: With R
%  Examples}. Springer, 4th edition. 
  
%     \bibitem[Jaynes(2003)]{Jaynes} Jaynes, E. T, (2003)
%  \textit{Probability theory: the logic of science}. Cambridge
%  University Press, 2003.
 
%  \bibitem[Sinai(1992)]{Sinai} Sinai, Y. G, (1992)
%  \textit{Probability theory: an introductory course}. Springer
%  Verlag, 1992.  

%\bibitem[{deGroot(1986)}]{DeGrootBlackwell}DeGroot, M. H. (1986)
%       A conversation with David Blackwell.
%\newblock \emph{Statistical Science}, \textbf{1}, 40--53.

%         \bibitem[Mosteller(1987)]{Mosteller} Mosteller, F, (1987)
%  \textit{Fifty challenging problems in probability with
%    solutions}. Courier Corporation, 1987.

%    \bibitem[MacKay(2003)]{MacKay} MacKay, D. J. C., (2003)
%  \textit{Information theory, inference, and learning
%  algorithms}. Cambridge University Press, 2003.

%\bibitem[{Lindley(2013)}]{Lindley}Lindley, D. V. (2013)
%   \textit{Understanding Uncertainty}. John Wiley and sons, 2013.


%\end{thebibliography}







\end{document}

%%% Local Variables:
%%% mode: latex
%%% TeX-master: t
%%% End:
